\documentclass[12pt]{article}%
\usepackage{amsfonts}
\usepackage{fancyhdr}
\usepackage{comment}
\usepackage[a4paper, top=2.5cm, bottom=2.5cm, left=2.2cm, right=2.2cm]%
{geometry}
\usepackage{times}
\usepackage{amsmath}
\usepackage{changepage}
\usepackage{amssymb}
\usepackage{graphicx}

\setcounter{MaxMatrixCols}{30}
\newtheorem{theorem}{Theorem}
\newtheorem{acknowledgement}[theorem]{Acknowledgement}
\newtheorem{algorithm}[theorem]{Algorithm}
\newtheorem{axiom}{Axiom}
\newtheorem{case}[theorem]{Case}
\newtheorem{claim}[theorem]{Claim}
\newtheorem{conclusion}[theorem]{Conclusion}
\newtheorem{condition}[theorem]{Condition}
\newtheorem{conjecture}[theorem]{Conjecture}
\newtheorem{corollary}[theorem]{Corollary}
\newtheorem{criterion}[theorem]{Criterion}
\newtheorem{definition}[theorem]{Definition}
\newtheorem{example}[theorem]{Example}
\newtheorem{exercise}[theorem]{Exercise}
\newtheorem{lemma}[theorem]{Lemma}
\newtheorem{notation}[theorem]{Notation}
\newtheorem{problem}[theorem]{Problem}
\newtheorem{proposition}[theorem]{Proposition}
\newtheorem{remark}[theorem]{Remark}
\newtheorem{solution}[theorem]{Solution}
\newtheorem{summary}[theorem]{Summary}
\newenvironment{proof}[1][Proof]{\textbf{#1.} }{\ \rule{0.5em}{0.5em}}

\newcommand{\Q}{\mathbb{Q}}
\newcommand{\R}{\mathbb{R}}
\newcommand{\C}{\mathbb{C}}
\newcommand{\Z}{\mathbb{Z}}

%%%%
% ME %
%%%%

\usepackage{listings}
\usepackage{color}
\usepackage{xcolor}
\usepackage{mdframed}

\definecolor{dkgreen}{rgb}{0,0.6,0}
\definecolor{gray}{rgb}{0.5,0.5,0.5}
\definecolor{mauve}{rgb}{0.58,0,0.82}

\lstset{%frame=tb,
language=Matlab,
aboveskip=3mm,
belowskip=3mm,
showstringspaces=false,
columns=flexible,
basicstyle={\small\ttfamily},
numbers=none,
numberstyle=\tiny\color{gray},
keywordstyle=\color{blue},
commentstyle=\color{dkgreen},
stringstyle=\color{mauve},
breaklines=true,
breakatwhitespace=true,
tabsize=3
}

\usepackage{outlines}
\usepackage{enumitem}
%\setenumerate[1]{label=\Roman*.}
%\setenumerate[2]{label=\Alph*.}
%\setenumerate[3]{label=\roman*.}
%\setenumerate[4]{label=\alph*.}

\usepackage{titlesec} 

\titleformat{\subsection}[runin]
  {\normalfont\large\bfseries}{\thesubsection}{1em}{}
\titleformat{\subsubsection}[runin]
  {\normalfont\normalsize\bfseries}{\thesubsubsection}{1em}{}

% for enumerate spacing
\newenvironment{tight_enum}{
\begin{enumerate}[label=\alph*.]
\setlength{\itemsep}{0pt}
\setlength{\parskip}{0pt}
}{\end{enumerate}}

\usepackage{exercise}

\setlength\parindent{0pt}

\newcommand{\code}[1]{\texttt{#1}}
\newcommand{\shp}{$>>>$ }
\newcommand{\tab}{\hspace*{45pt}}

\usepackage{multicol}
\setlength{\columnsep}{0.6cm}

\begin{document}

\title{Homework 09: \\ NumPY and Matplotlib}

\author{EEC 3150}
\date{}

\maketitle

\begin{center}
\fbox{\fbox{\parbox{5.5in}{\centering
For each Python exercise, write a separate Python script. Submit all scripts on Blackboard under the appropriate assignment. Name your script files with the following format: \\
\vspace{10pt}
HW$<$assignment number$>$\_Ex$<$exercise number$>$\_$<$FirstInitial$><$LastName$>$.py \\
\vspace{10pt}
Example: HW01\_Ex01\_GWest.py \\
\vspace{10pt}
For short answer or math exercises, either do your work on a sheet of paper and upload a picture of it to Blackboard or submit some kind of text file or Word document to Blackboard. You can place all such exercises in a single document as long as you clearly label each exercise. These files should follow a similar naming convention as the Python scripts.
}}}
\end{center}


\pagebreak


\begin{multicols*}{2}



%%% NEW EXERCISE %%%
\begin{Exercise}[origin={10 points, Python}]

Create and print the following array WITHOUT directly coding any of the elements (only use numpy functions  or operations): \\

\code{
\begin{tabular}{ r c c l }
[[ 1 & 3 & 5 & 7  ] \\[0pt]
 [ 9 & 11 & 13 & 15 ] \\[0pt]
 [17 & 19 & 21 & 23 ] \\[0pt]
 [25 & 27 & 29 & 31 ]]
\end{tabular}
}

\end{Exercise}


%%%% NEW EXERCISE %%%
%\begin{Exercise}[origin={10 points, Python}]
%
%Create and print this 2D array of complex numbers WITHOUT directly coding any of the elements (only use numpy functions or operations): \\
%
%\code{
%\begin{tabular}{ r c l }
%[[ 1.+1.j & 2.+1.j & 3.+1.j ] \\[0pt]
% [ 1.+2.j & 2.+2.j & 3.+2.j ] \\[0pt]
% [ 1.+3.j & 2.+3.j & 3.+3.j ]]
%\end{tabular}
%}
%
%\end{Exercise}


%%% NEW EXERCISE %%%
\begin{Exercise}[origin={10 points, Python}]

\begin{enumerate}[leftmargin=*]
	\item Write a function \code{Boundary(n,value)} which will create an $n\times n$ matrix in which the elements on the top, bottom, left, and right boundaries are equal to \code{value} while the inside elements are equal to 0.
	\item Demonstrate your function by calling it and printing the output
\end{enumerate}

Here is an example: \code{Boundary(4,1)} should produce: \\

\code{
\begin{tabular}{ r c c l }
[[ 1 & 1 & 1 & 1 ] \\[0pt]
 [ 1 & 0 & 0 & 1 ] \\[0pt]
 [ 1 & 0 & 0 & 1 ] \\[0pt]
 [ 1 & 1 & 1 & 1 ]]
\end{tabular}
}

\end{Exercise}


%%% NEW EXERCISE %%%
\begin{Exercise}[origin={10 points, Python}]

\begin{enumerate}[leftmargin=*]
	\item Write a function \code{Truncate(arr)} which takes an array \code{arr} of integers \underline{\textbf{of any shape}} and returns a new array with the following properties:
	\begin{itemize}
		\item Set all values where \code{arr} is even to 0
		\item Set all values where \code{arr} is odd to 1
	\end{itemize}
	\item Demonstrate your function by generating an array of random integers on calling your function on it. Print the output.
\end{enumerate}

Here is an example: \\ \code{Truncate(np.array([-1,2,11,-4]))} should produce: \\

\code{[1 0 1 0]}

\end{Exercise}


%%% NEW EXERCISE %%%
\begin{Exercise}[origin={20 points, Python}]

Plot a function of your choice on a domain of your choice.

\begin{enumerate}[leftmargin=*]
	\item Use \code{np.linspace} to create an array of $x$ values (you choose the parameters)
	\item Create an array of $y$ values by evaluting some function of your choice on the array of $x$ values
	\item Plot attributes
	\begin{itemize}
		\item Plot the $x,y$ data in two ways on the same figure: 1) a green dashed line and 2) red dots
		\item Set the makersize and linewidth to something other than the default
		\item Set xlim, ylim to some reasonable values to display the entire plot
	\end{itemize}
\end{enumerate}

\end{Exercise}


%%% NEW EXERCISE %%%
\begin{Exercise}[origin={20 points, Python}]

The exercise will have you generate random data, plot it, and index it.

\begin{enumerate}[leftmargin=*]
	\item Generate random data
	\begin{itemize}[leftmargin=*]
		\item Create 1D \code{x,y} arrays of 1000 random numbers between -2 and 2
	\end{itemize}
	\item Plot the data
	\begin{itemize}[leftmargin=*]
		\item Plot the \code{x,y} arrays as blue dots
	\end{itemize}
	\item Index the data
	\begin{itemize}[leftmargin=*]
		\item Through whatever means you choose, create \code{x2,y2} arrays which only have the points where $-1.5<x<1$ and $-1<y<0.9$
	\end{itemize}
	\item Plot the new data
	\begin{itemize}[leftmargin=*]
		\item Plot the \code{x2,y2} arrays on the same plot as red x's with an increased \code{markersize} for readability
	\end{itemize}
\end{enumerate}

\end{Exercise}

More problems on next page...

\vfill\null
\pagebreak

%%% NEW EXERCISE %%%
\begin{Exercise}[origin={20 points, Python}]

\begin{enumerate}[leftmargin=*]
	\item Create a function to generate a 2D rotation matrix (2D array) (NOTE: sin and cos take radians). For reference, here is the matrix:
	\begin{equation*}
		\begin{pmatrix}
			\cos(\theta) & -\sin(\theta) \\
			\sin(\theta) & \cos(\theta)
		\end{pmatrix}
	\end{equation*}
	\item Create a rotation matrix \code{R} with $\theta=2\pi/n$ where $n$ is an integer variable of your choice (it should be $10\le n \le 100$)
	\item Create the vector (1D array) \code{x = [1 0]}
	\item Create an array of vectors (2D array) by iteratively applying the rotation matrix to the vector \code{x}. This should produce the points on the unit circle
	\item Plot the points
\end{enumerate}

\end{Exercise}











%\vfill\null
%\columnbreak

\vfill\null

\end{multicols*}

\end{document}






